
% AI-assisted factual update 2026-11-27: pending supervisor permission

This appendix defines the terms used in the thesis. Terms of the case study come first; terms that only the earlier demonstrator uses are in the second table.

\section{Case-study terms}

{\small
\begin{longtable}{>{\raggedright\arraybackslash}p{4.63cm}>{\raggedright\arraybackslash}p{9.57cm}}
\toprule
Term & Definition \\
\midrule
\endhead
\textbf{Agreement with COS} & How far the candidate-change flags coincide with land-use change recorded in the COS map, given as precision, recall and F1. It compares two maps; it is not accuracy against field observations. \\
\textbf{Band} & One wavelength interval recorded by a sensor, for example red (Sentinel-2 B04) or near infrared (B08). \\
\textbf{Bootstrap (spatial block)} & Repeated resampling of whole blocks of pixels or whole COS polygons, with replacement, to estimate a range for precision and recall. Blocks are used because neighbouring pixels are not independent. \\
\textbf{Candidate change} & A pixel whose index difference exceeds the threshold and that lies in a connected patch of at least 1 ha. It says where the surface changed, not why. \\
\textbf{Cloud masking} & Removing pixels covered by cloud, cloud shadow or other unusable conditions. The case study keeps only SCL classes 4 (vegetation), 5 (not vegetated) and 6 (water). \\
\textbf{Control pair} & Two images of the same place taken 20 days apart in the same summer. Their difference estimates how much the index varies without land-use change. \\
\textbf{Data status} & The label every published number carries: public observation, observation-derived result, unverified interpretation or synthetic demonstration. \\
\textbf{Index difference ($\Delta$)} & The index value on the later date minus the value on the earlier date, per pixel. \\
\textbf{Manifest} & A JSON file written with each run that records the input item ids, code commit, software versions, settings and a SHA-256 hash of every output. \\
\textbf{Patch} & A group of flagged pixels that touch each other, including diagonally (8-connected). \\
\textbf{Processing baseline} & The version of the ESA processing software that produced a Sentinel-2 product, for example 05.00 or 05.11. Products from baseline 04.00 onward carry a reflectance offset. \\
\textbf{Surface reflectance} & The share of sunlight reflected by the ground after correction for the atmosphere; the content of a Sentinel-2 Level-2A product. \\
\textbf{Threshold rule} & The case-study rule for choosing the threshold before looking at the change map: the smallest value on the configured grid that is at or above the 99th percentile of the control pair's absolute difference. \\
\textbf{Tile} & One 100 $\times$ 100 km square of the MGRS grid in which Sentinel-2 products are delivered, for example 29SNC. \\
\bottomrule
\end{longtable}
}

\section{Earlier demonstrator terms}

{\small
\begin{longtable}{>{\raggedright\arraybackslash}p{3.81cm}>{\raggedright\arraybackslash}p{10.39cm}}
\toprule
Term & Definition \\
\midrule
\endhead
\textbf{Backscatter} & The part of a radar signal reflected back to the antenna (Sentinel-1). \\
\textbf{Contamination rate} & The share of observations that an Isolation Forest is told to treat as outliers. \\
\textbf{Isolation Forest} & An unsupervised method that scores how easily an observation is separated from the rest by random splits. Used only in the earlier demonstrator, on synthetic data. \\
\textbf{Radiance} & Light energy leaving a surface per unit area, angle and wavelength; night-time light products report it in nanowatts per square centimetre per steradian. \\
\textbf{SWIR} & Short-wave infrared, about 1.4 to 3.0 micrometres. Sentinel-2 band B11 (about 1.6 micrometres) is used for BSI. \\
\textbf{Z-score} & The number of standard deviations an observation lies from the mean of a reference period. \\
\bottomrule
\end{longtable}
}
